\chapter{Реализация базового алгоритма BC}
\label{ch:chap2}

В соответствии с техническим заданием в качестве базового решения требовалось реализовать визуомоторную политику клонирования поведения и обеспечить успешность не менее 75\,\% эпизодов на наборе данных robomimic~v0.1. В исходной формулировке задания базовый алгоритм обозначался как BC-RNN, однако по итогам анализа литературы (см.~Главу~\ref{ch:chap1}) рекуррентная часть заменена на каузальный трансформер декодерного типа: он точнее моделирует длинные временные зависимости в окне наблюдений и не наследует ограничений рекуррентных сетей, связанных с затуханием градиента и последовательной обработкой. Для извлечения признаков изображения вместо обучаемой с нуля сверточной сети выбран самообучаемый ViT DINOv2 \cite{oquab2024dinov2}: он обеспечивает более информативные и устойчивые представления при ограниченном объеме демонстраций и является отправной точкой для последующего перехода к квантованию визуальных представлений, рассматриваемому в Главе~\ref{ch:chap3}.

\section{Архитектура политики}
\label{sec:bc-architecture}

Реализация выполнена на языке Python с использованием PyTorch и поверх библиотеки robomimic \cite{mandlekar2021robomimic}, что позволило переиспользовать стандартную инфраструктуру загрузки демонстраций, аугментации, обучения и оценки политики. На каждом временном шаге $t$ политика получает наблюдение $o_t = (s_t, I_t^{\text{eih}}, I_t^{\text{agent}})$, где $s_t$ -- низкоразмерное состояние робота (положение и ориентация захвата $\text{robot0\_eef\_pos}$, $\text{robot0\_eef\_quat}$ и угол схвата $\text{robot0\_gripper\_qpos}$), а $I_t^{\text{eih}}$ и $I_t^{\text{agent}}$ -- кадры с камеры, закрепленной на захвате (eye-in-hand), и сценической камеры. На основании окна последних $T = 10$ наблюдений политика возвращает действие $a_t$ -- 7-мерный вектор приращения положения и ориентации захвата вместе с командой схвату.

Политика построена по типовой для визуомоторного клонирования поведения схеме, состоящей из трех частей: (а)~общего энкодера наблюдения, формирующего на каждом шаге компактное представление $h_t$; (б)~авторегрессивного трансформера, преобразующего последовательность $h_{t-T+1:t}$ в скрытое состояние политики; и (в)~стохастической головы, выдающей распределение действий $\pi_\theta(a_t \mid h_{t-T+1:t})$.

\subsection{Зрительный энкодер на базе DINOv2}
\label{subsec:bc-encoder}

Для каждой из двух камер входное изображение $I_t \in \mathbb{R}^{3\times H \times W}$ независимо обрабатывается одним и тем же зрительным энкодером. Перед подачей на сеть применяется случайная кадрировка (RandomCrop размером $76{\times}76$, в режиме оценки -- центральная). Кадрировка используется как штатная аугментация: согласно \cite{mandlekar2021robomimic}, наличие случайных пиксельных сдвигов заметно повышает устойчивость визуомоторных политик, тогда как их отключение ухудшает успешность захвата.

В качестве базовой зрительной сети выбран ViT DINOv2 размера S/14, импортируемый напрямую из официального репозитория модели. Изображение приводится к разрешению $224{\times}224$, что при размере патча $14{\times}14$ дает сетку из $16{\times}16 = 256$ патч-токенов. Каждый токен -- вектор размерности 384. Чтобы получить из последовательности патч-токенов карту признаков, удобную для последующей обработки, токены классификации и регистровые токены отбрасываются, а оставшиеся патч-токены переставляются в тензор формы $[B, 384, 16, 16]$.

Полное обучение DINOv2 нецелесообразно: модель содержит порядка 22\,млн параметров, обучена на 142~миллионах изображений и при дообучении на нескольких сотнях демонстраций склонна терять часть заложенной общности. В то же время полное замораживание весов также не оптимально: предобученные представления получены на естественных изображениях, тогда как кадры из MuJoCo имеют свои текстурные и колориметрические особенности. В качестве компромисса принят режим частичного дообучения: все блоки трансформера, кроме последнего, а также эмбеддинг патчей и позиционные эмбеддинги, оставлены замороженными; последний блок и финальная нормализация обучаются вместе с остальной сетью. Такая схема позволяет адаптировать высокоуровневые признаки к симулированной среде, не разрушая при этом низкоуровневые свойства предобученного энкодера. Для уменьшения потребления памяти прямой проход по замороженной части выполняется в режиме \texttt{torch.no\_grad()}.

Поверх карты $[B, 384, 16, 16]$ в базовой конфигурации применяется обучаемый слой SpatialSoftmax с 32 ключевыми точками. Этот модуль преобразует пространственную карту в набор $(x, y)$-координат точек максимума активации и широко применяется в визуомоторных политиках \cite{mandlekar2021robomimic}. Получаемое представление пропускается через линейный слой размерности 64 и конкатенируется с признаками второй камеры и низкоразмерным состоянием~$s_t$, формируя итоговое наблюдение шага $h_t$.

\subsection{Генератор действий на базе трансформера}
\label{subsec:bc-policy}

Сформированная последовательность $h_{t-T+1:t}$ подается в декодерный трансформер с каузальной маской, обеспечивающей авторегрессивное предсказание по схеме, аналогичной языковым моделям с каузальным самовниманием \cite{vaswani2017attention, radford2018improving}. Используется 6 слоев с эмбеддингом размерности~512 и 8 головами внимания, функция активации в feed-forward блоках -- GELU; вероятности прореживания (на эмбеддинги, веса внимания и выход блока) равны 0{,}1. Позиционные эмбеддинги задаются обучаемым параметром по числу временных шагов окна, что согласуется с фиксированной длиной контекста~$T$. Для устойчивости обучения на каждый шаг $t$ независимо вычисляется потеря (опция \texttt{supervise\_all\_steps}), что эквивалентно стандартной для языковых моделей причинно-маскированной NLL по всей последовательности.

Стохастическая голова реализует мультимодальную политику в виде смеси гауссиан (GMM) с пятью компонентами; среднее каждой компоненты получается линейной проекцией скрытого состояния трансформера, стандартное отклонение -- через активацию softplus с нижним порогом $10^{-4}$. Такой выбор обусловлен наблюдением \cite{mandlekar2021robomimic} о мультимодальности человеческих демонстраций: усреднение по режимам приводит к политике, которая ни одному из них не следует. На этапе оценки используется режим пониженного шума (\texttt{low\_noise\_eval}), что эквивалентно выбору моды распределения и улучшает воспроизводимость эпизодов развертывания.

Целевая функция при обучении -- отрицательная логарифмическая правдоподобность действия эксперта $a_t^{\ast}$ под политикой:
\begin{equation}
    \label{eq:bc-loss}
    \mathcal{L}_{\text{policy}}(\theta) = -\mathbb{E}_{(o_{t-T+1:t}, a_t^{\ast}) \sim \mathcal{D}}
    \log \pi_\theta\!\left(a_t^{\ast} \mid o_{t-T+1:t}\right),
\end{equation}
где $\mathcal{D}$ -- набор демонстраций. Оптимизация ведется методом Adam с начальной скоростью $10^{-4}$ без L2-регуляризации; норма градиента ограничивается значением 100.

\section{Обучение и условия эксперимента}
\label{sec:bc-training}

Обучение выполнено на наборе robomimic~v0.1 для задачи Lift в среде robosuite (MuJoCo). Используется только подмножество демонстраций категории \emph{Proficient-Human} (\texttt{ph}), собранных квалифицированным оператором. Из наблюдений намеренно исключены истинные позы объекта (\texttt{object} в терминах robomimic): это смещает решение от аналитического в сторону визуомоторного, что соответствует постановке задачи. Длина окна наблюдений совпадает с длиной контекста трансформера и равна 10 шагам. Размер мини-батча -- 16, число эпох -- 100, эпохой считаются 100 шагов градиентного спуска. Каждые 10 эпох проводится оценка политики путем развертывания в исходной среде: запускается 10 тестовых эпизодов с горизонтом 400 шагов; эпизод считается успешным, если задача выполнена до истечения горизонта.

\section{Результаты}
\label{sec:bc-results}

При тестировании описанной выше политики на задаче Lift из robomimic~v0.1 был достигнут уровень успешности $\text{Success Rate} = 0{,}95$, что превышает требуемый порог в 0{,}75 и подтверждает работоспособность базовой реализации. Полученный результат используется в дальнейшем как точка сравнения для алгоритма с квантованием визуальных представлений, описанного в Главе~\ref{ch:chap3}.

\endinput
